%% einfart minimal article starter — the plain face of Jinwen XU's
%% "minimalist" class series (CTAN: einfart).
%%
%% The einfart class itself is NOT included in this download: it ships with
%% every modern TeX distribution (and with latextemplates.com's compiler),
%% so bundling a copy would only risk pinning you to an older version than
%% your own system already has.
%%
%% Class documentation: texdoc einfart   |   https://ctan.org/pkg/einfart
%% Licence: LPPL-1.3c, (c) 2021-2026 Jinwen XU -- verified directly in the
%% shipped einfart.cls header.
%%
%% einfart is the PLAIN counterpart to the already-catalogued "Stylish
%% Article" substitute (slug stylish-article), which uses lebhart -- the
%% COLOURFUL face of the same class series (same author, same LPPL-1.3c).
%% Where lebhart adds coloured section rules and theorem boxes, einfart is
%% deliberately unstyled: a clean, minimal article face with careful
%% typography and no colour.
%%
%% ENGINE: einfart requires XeLaTeX or LuaLaTeX (fontspec-based); see
%% engine.txt. It also loads ctex unconditionally for CJK support -- real
%% Chinese text renders correctly under this class (verified in the
%% preview), resolved via kpathsea/Fandol, not tofu.
\documentclass{einfart}

\title{Your Title Here}
\author{Ada Lovelace}
\date{\today}

\begin{document}

\maketitle

\begin{abstract}
Replace this with your abstract. A good abstract states the problem, what
you did, and what you found --- in that order, and without citations.
\end{abstract}

\section{Introduction}
Start here. Cite things like this~\cite{knuth1984texbook}. This class's
default typography uses TeX Gyre Pagella, a Palatino-like serif with wide
Unicode coverage, resolved directly by fontspec rather than a bare
fontconfig lookup.

\section{CJK support}
\label{sec:cjk}
einfart loads \texttt{ctex} unconditionally, so Chinese text can appear
directly inside an otherwise English document without any extra setup:
这是一段中文示例文字，用来验证中日韩文字能够正确渲染。

\section{Mathematics}
\label{sec:math}
Standard math environments work as usual:
\begin{equation}
  \label{eq:main}
  e^{i\pi} + 1 = 0.
\end{equation}
Equation~\eqref{eq:main} is referenced like this.

\section{Conclusion}
Say what it means.

\bibliographystyle{plain}
\bibliography{references}

\end{document}
